\section{Physical coordinates of the two local outputs}
\label{sec:peps_local_output_coordinates}

We use the actual physical and discarded register lists in the regional
decomposition of \cite[proof of Theorem~5.2]{OpenAI2026PolynomialPEPS}.
% Source: 04-compression.tex, lines 409--480, at
% adc7f1241b42e322a6451854ab7e4b4c146bf78a.
All coordinate bases in this
section are finite orthonormal bases. The private dimensions are unrestricted.

\begin{definition}[Bases adapted to the party partition]
    \label{def:peps_partition_output_bases}
    \lean{TNLean.PEPS.PairEffect.Layout.partitionBasis}
    \lean{TNLean.PEPS.PairEffect.Layout.partitionAppendBasis}
    \lean{TNLean.PEPS.PairEffect.Layout.partitionOutputIso}
    \lean{TNLean.PEPS.PairEffect.Layout.regionalOutputIso}
    \leanok
    \uses{def:peps_party_partition}
    Write the output as a physical register list followed by a discarded
    register list. Partition each list into the affected and exterior owners.
    Choose orthonormal bases $(a_x)$ and $(e_y)$ of the two physical memories,
    and $(d_i)$ and $(f_j)$ of the two discarded memories. Transport their
    tensor products by the canonical register permutations. This gives
    a basis with coordinates $((x,y),(i,j))$ in physical/discard order,
    and a basis with coordinates $((x,i),(y,j))$ in affected/exterior order.

    On the affected memory, the output isometry sends
    $a_x\otimes d_i$ to $e_x\otimes d_i$, where $(e_x)$ is the standard
    coordinate basis. The exterior isometry is defined in the same way.
    If a register list is identified with its physical/discard concatenation,
    these maps include the canonical identification of equal lists.
\end{definition}

\begin{theorem}[Regrouping physical and discarded coordinates]
    \label{thm:peps_partition_output_coordinates}
    \lean{TNLean.PEPS.PairEffect.Layout.memCongr_tmul_heq}
    \lean{TNLean.PEPS.PairEffect.Layout.partitionIso_partitionBasis_tmul}
    \lean{TNLean.PEPS.PairEffect.Layout.partitionBasis_coordinates}
    \lean{TNLean.PEPS.PairEffect.Layout.partitionOutputIso_basis}
    \lean{TNLean.PEPS.PairEffect.Layout.partitionOutputIso_coordinates}
    \lean{TNLean.PEPS.PairEffect.Layout.partitionBasis_coordinates_of_eq}
    \lean{TNLean.PEPS.PairEffect.Word.repr_effectMap}
    \leanok
    \uses{def:peps_partition_output_bases}
    For every output vector $z$, its coordinate at $((x,y),(i,j))$
    in physical/discard order is its coordinate at $((x,i),(y,j))$
    after the affected/exterior partition. The affected output isometry
    preserves the pair $(x,i)$ of coordinates, and similarly on the exterior.
    For $v\in\C^X\otimes D$, taking the partial inner product with $e_x$
    and then the $d_i$ coordinate gives exactly the $(x,i)$ coordinate of $v$.
\end{theorem}

\begin{proof}\leanok
    \uses{def:peps_partition_output_bases}
    On an elementary basis tensor the two orders are
    $(a_x\otimes e_y)\otimes(d_i\otimes f_j)$ and
    $(a_x\otimes d_i)\otimes(e_y\otimes f_j)$.
    The canonical partition sends the first to the second. Taking inner
    products with an arbitrary $z$ gives the coordinate identity because
    the partition is isometric. Partial inner products give the stated
    coordinate projection first on elementary tensors and then by linearity.
\end{proof}

\begin{theorem}[Output coordinates under owner relabelling]
    \label{thm:peps_owner_output_coordinates}
    \lean{TNLean.PEPS.PairEffect.Layout.mapOwner_tensorBasis_coordinates}
    \lean{TNLean.PEPS.PairEffect.Layout.mapOwner_mapOwner_tensorBasis_coordinates}
    \lean{TNLean.PEPS.PairEffect.Layout.mapOwner_mapOwner_partitionBasis_coordinates}
    \leanok
    \uses{thm:peps_partition_output_coordinates}
    Relabel owners by any map, or successively by two maps. Transport the
    physical and discarded bases by the corresponding memory isometries.
    Every physical/discard coordinate of the original vector equals the
    same coordinate of the relabelled vector. In particular, partitioning
    after the two relabellings gives the coordinate order
    $((x,i),(y,j))$ described above. No injectivity of the owner maps is needed.
\end{theorem}

\begin{proof}\leanok
    \uses{thm:peps_partition_output_coordinates}
    Owner relabelling preserves the ordered register spaces and commutes
    with concatenation. Its isometry therefore sends each transported
    physical/discard basis tensor to the original tensor. Preserve inner
    products, apply this twice, and then apply the partition identity.
\end{proof}

\begin{theorem}[Coordinates of a local Schmidt expansion]
    \label{thm:peps_local_schmidt_output_coordinates}
    \lean{TNLean.PEPS.PairEffect.Layout.partition_coordinates_eq_separated}
    \lean{TNLean.PEPS.PairEffect.Layout.mapOwner_mapOwner_coordinates_eq_separated}
    \leanok
    \uses{thm:peps_owner_output_coordinates}
    Let $V,W$ be the two actual local output maps, let $F,G$ be input
    isometries, and suppose that the affected/exterior partition of an
    output vector $z_k$ is
    \begin{align}
        \sum_t\sqrt{\tau_t}\,
        VF e_{(k,t)}\otimes WG e_t.
        \label{eq:peps_local_schmidt_output_vector}
    \end{align}
    In the partition-adapted physical/discard bases, the coordinate of $z_k$
    at $((x,y),(i,j))$ is
    \begin{align}
        \sum_t\sqrt{\tau_t}\,
        \langle d_i,C_{x,V}F e_{(k,t)}\rangle
        \langle f_j,C_{y,W}G e_t\rangle,
        \label{eq:peps_local_schmidt_output_coordinate}
    \end{align}
    where $C_{x,V}$ and $C_{y,W}$ are the physical-coordinate projections
    of the two local outputs. The same identity holds for the original
    output before either of the two owner relabellings, in the transported
    bases. This identity does not require contractivity or normalization
    in addition to the displayed vector equality.
\end{theorem}

\begin{proof}\leanok
    \uses{thm:peps_partition_output_coordinates,thm:peps_owner_output_coordinates}
    Substitute (\ref{eq:peps_local_schmidt_output_vector}) into the partition
    coordinate identity. Linearity gives the sum over $t$, and the tensor
    inner product gives the product of the two discarded coordinates.
    The partial-inner-product identity identifies these coordinates with
    $C_{x,V}$ and $C_{y,W}$. Transport the bases back through the two owner maps.
\end{proof}

\begin{theorem}[Original physical dimensions after the two owner maps]
    \label{thm:peps_affected_original_dimensions}
    \lean{TNLean.PEPS.PairEffect.Layout.restrict_mapOwner}
    \lean{TNLean.PEPS.PairEffect.Layout.restrictMapOwnerIso}
    \lean{TNLean.PEPS.PairEffect.Layout.affectedOutputIso}
    \lean{TNLean.PEPS.PairEffect.affectedFamilyPhysicalBasis}
    \lean{TNLean.PEPS.PairEffect.finrank_affectedFamilyPhysicalLayout}
    \leanok
    \uses{thm:peps_family_regional_dimensions}
    Let the original parties form a finite list without repetitions,
    containing every party. Party $p$ has physical space $\C^{d_p}$.
    Relabel affected parties individually, identify all exterior parties,
    and then map owners to the two sides. The selected physical memory
    is canonically isometric to the original memory on that side.
    For an affected set $A$, its actual orthonormal coordinates are indexed
    by assignments $s_p\in\{0,\ldots,d_p-1\}$ for $p\in A$, and its
    dimension is $\prod_{p\in A}d_p$.
\end{theorem}

\begin{proof}\leanok
    \uses{thm:peps_family_regional_dimensions}
    Selecting after relabelling is the same as relabelling the registers
    selected by the composite owner test. For the two stated owner maps,
    that test is exactly membership in $A$. Transport the original regional
    tensor basis through these isometries and count its assignments.
\end{proof}
